\chapter{Aspectos Computacionales de FEniCSx}
\label{app:fenics}

Este apéndice complementa la exposición práctica de los capítulos anteriores con una discusión técnica de los algoritmos numéricos que operan bajo el capó de FEniCSx. Comprender estas elecciones algorítmicas permite al ingeniero escalar sus simulaciones desde prototipos 2D hasta modelos industriales de alta fidelidad.

%==================================================================================
\section{Solvers Directos vs.~Iterativos}
%==================================================================================

La resolución del sistema lineal $A\bm{x} = \bm{b}$ resultante de la discretización por elementos finitos constituye el cuello de botella computacional en la mayoría de las simulaciones. La elección del solver depende fundamentalmente de la escala del problema.

\textbf{Solvers directos:} Para problemas 2D de tamaño medio (hasta $10^5$--$10^6$ grados de libertad), los métodos directos basados en factorización LU son robustos y predecibles. En FEniCSx, la opción por defecto \texttt{"pc\_type": "lu"} invoca al solver MUMPS (\textit{MUltifrontal Massively Parallel Solver}), que factoriza la matriz global en una descomposición triangular y resuelve los sistemes por sustitución. La principal ventaja es la independencia del número de condición: MUMPS converge siempre que la matriz sea invertible.

\textbf{Solvers iterativos:} En problemas 3D industriales, la matriz de rigidez es notablemente más grande y dispersa. Los métodos directos sufren un crecimiento exponencial del coste de factorización ($\mathcal{O}(N^2)$ en 3D). Aquí, los métodos iterativos como el de \textbf{Gradientes Conjugados (CG)} combinados con un precondicionador de multigrid algebraico (AMG) resultan escalables. El precondicionador AMG aproxima la inversa de $A$ mediante una jerarquía de operadores en mallas de distinta resolución, acelerando drásticamente la convergencia.

FEniCSx delega esta elección al backend PETSc, que expone los parámetros \texttt{ksp\_type} (tipo de solver iterativo) y \texttt{pc\_type} (tipo de precondicionador). Un configuración típica para problemas 3D grandes es:
\begin{codigo}[language=Python]
petsc_options={"ksp_type": "cg", "pc_type": "hypre", "pc_hypre_type": "boomeramg"}
\end{codigo}

%==================================================================================
\section{Linealización de Newton-Raphson}
\label{sec:newton-raphson}
%==================================================================================

Cuando la forma débil es no lineal en la incógnita $u$, la ecuación residual $F(u) = 0$ debe resolverse mediante un proceso iterativo. El método de \textbf{Newton-Raphson} es el algoritmo de referencia, y constituye el motor subyacente de los solvers no lineales de FEniCSx.

Dado un problema genérico no lineal $F(u) = 0$, donde $F: V \to V'$ es un operador diferenciable, el algoritmo procede como sigue:

\begin{enumerate}
    \item \textbf{Inicialización:} Elegir una aproximación inicial $u_0 \in V$.
    \item \textbf{Iteración:} Para $k = 0, 1, 2, \dots$ hasta convergencia:
    \begin{enumerate}
        \item Calcular la \textbf{matriz jacobiana} (derivada de Fréchet) $J_F(u_k)$ del operador $F$ evaluada en $u_k$.
        \item Resolver el sistema lineal de corrección:
        \begin{equation}
            J_F(u_k) \, \delta u = -F(u_k).
            \label{eq:newton-step}
        \end{equation}
        \item Actualizar la solución:
        \begin{equation}
            u_{k+1} = u_k + \delta u.
            \label{eq:newton-update}
        \end{equation}
    \end{enumerate}
    \item \textbf{Criterio de convergencia:} Detener cuando la norma de la corrección sea menor que una tolerancia prescrita:
    \begin{equation}
        \norm{\delta u}_V < \varepsilon,
        \label{eq:newton-tol}
    \end{equation}
    o bien cuando la norma del residuo relativo $\norm{F(u_k)} / \norm{F(u_0)}$ descienda por debajo de un umbral.
\end{enumerate}

La convergencia del método de Newton es \textbf{cuadrática} en una vecindad suficientemente cercana a la solución, siempre que $J_F$ sea Lipschitz-continua y no singular en la raíz. En la práctica, para problemas de multifísica con fuertes no linealidades (cambio de fase, convección dominante, etc.), puede ser necesario recurrir a estrategias de globalización como la \textbf{line search} o el método de \textbf{relajación} (damping) para garantizar la convergencia desde condiciones iniciales arbitrarias.

\begin{notatecnica}
En FEniCSx, el solver no lineal \texttt{NewtonSolver} automatiza todo este ciclo: el usuario proporciona únicamente la forma variacional residual $F(u; v)$ y el jacobiano $J = \mathrm{D}F$ se puede derivar automáticamente mediante diferenciación simbólica de UFL. Este es el motor que se utilizará en la Parte~II del libro para abordar problemas termoelásticos y de flujo no lineal.
\end{notatecnica}
